\section{Điểm nổi bật và đóng góp của dự án}
    \subsection{Kết nối thị giác máy tính với chuỗi thời gian và xử lý dữ liệu khởi động}
Dự án chuyển đổi số đếm phương tiện có timestamp thành chuỗi đo theo phút, dùng chung định nghĩa đặc trưng khi huấn luyện và chạy trực tiếp để tạo ra đường dự báo và ETA. Mọi bước đều có đầu vào và đầu ra kiểm tra được thay vì chỉ trình diễn một mô hình riêng lẻ.

Để xử lý việc mô hình cần 61 điểm phút ngay khi khởi động, hệ thống dùng chuỗi tham chiếu học từ data1 qua một nhánh tạm thời và thay thế dần khi tích lũy đủ quan sát thật. Dữ liệu tham chiếu không được đưa vào estimator như dữ liệu đo thật, và việc mất camera sau khởi động cũng không bị che bằng đường nội suy.

\subsection{Vận hành bất đồng bộ, cố định nguồn camera và lưu trữ đánh giá}
Giao diện cho phép chọn camera hiển thị, nhưng nguồn dự báo luôn cố định là Camera 1 để không làm trộn lẫn lịch sử chuỗi thời gian giữa các vị trí.

Tác vụ thu ảnh/YOLO và dự báo được đưa ra khỏi luồng giao diện; mô hình chỉ tải một lần và biểu đồ tái sử dụng đối tượng vẽ nhằm giữ độ phản hồi cho giao diện. Cơ sở dữ liệu SQLite lưu vết đầy đủ timestamp, origin, thời điểm phát, giá trị thô/làm mượt và model ID. Khi đánh giá, hệ thống chỉ lấy lần phát đầu tiên của mỗi origin và chỉ chấm điểm khi có đủ dữ liệu tương lai với độ phủ yêu cầu, tránh việc chọn lại phiên bản dự báo thuận lợi sau khi đã biết kết quả.

\subsection{Kiểm tra tương thích mô hình và kiểm thử hướng lỗi thực tế}
Hàm \texttt{validate\_bundle()} kiểm tra cấu trúc bundle trước khi suy luận, bao gồm yêu cầu đủ 60 estimator và đúng thứ tự cột đặc trưng. Khi huấn luyện mới, script chuyển sang lưu dưới tên candidate để tải lại và đối chiếu dự báo trước khi lưu, không tự ý ghi đè lên artifact đang vận hành.

Các trường hợp kiểm thử bao phủ trực tiếp các lỗi vận hành thực tế: đặc trưng không sử dụng tương lai, đúng thứ tự lag 60, lệch schema, thiếu model, timestamp dự báo, quy tắc 3/5, mất mẫu, phút thưa, đổi camera, khôi phục SQLite, giữ phiên bản dự báo và ghi nhận khi có lỗi nhận diện.   